\documentclass[main.tex]{subfiles}


\begin{document}

% \AddToShipoutPictureFG{%
% \begin{tikzpicture}[overlay,remember picture]
% \path (current page.south west) -- (current page.north east)
% node[midway,scale=1,color=gray,sloped,opacity=0.4] {\textnormal{Кравченко Игорь Игоревич\hspace{1cm} +79010144910\hspace{1cm} super.antig@yandex.ru}};
% \end{tikzpicture}
% }


% %from pagenumber to up
% \phantomsection 
%\hypertarget{toc}{}

 

 

% номер секции вручную
\setcounter{section}{11
}
\addtocounter{section}{-1} 

\section{Движение по окружности}


Для начала следует рассмотреть движение по окружности или ее части~--- дуге, происходящее с \emph{постоянной величиной скорости}.

Пусть имеется колесо от тележки (шину сняли), ось которого закрепили горизонтально. В некоторой точке на краю колеса оказался муравей (тело); колесо привели в равномерное вращение. На рис.\,\ref{fig:p21} показан опыт с движением этого муравья (положения тела фиксируют ежесекундно).

\begin{figure}[H]
\centering

    \begin{tikzpicture}[scale=.7,font=\small,            line cap=round,                             >={Stealth[inset=0pt,round,length=3mm, angle'=20]},vec/.style={>={Stealth[inset=0pt,round,length=3.5mm, angle'=20]}}]

\filldraw[draw=gray,fill=lightgray,semitransparent] (0,0) circle (3);
\draw[densely dashed] (0,3) -- (0,0) node[midway,left,black] {$R$};
\node at (0,0) [circle,fill=gray,inner sep=0pt,minimum size=1mm,label=below:$O$] {};


%vec a
\draw[->,red,thick,vec] (0,3) --++ (-90:1) node[black,right] {$a_\text{ц}$};
\draw[->,red,thick,vec] (45:3) --++ (225:1) node[black,below right,inner xsep=0] {$a_\text{ц}$};
\draw[->,red,thick,vec] (0:3) --++ (180:1) node[black,below] {$a_\text{ц}$};


%vec v
\draw[->,NavyBlue,thick,vec] (0,3) --++ (0:1.5) node[black,above] {$v_\text{к}$};
\node at (0,3) [circle,fill=,inner sep=0pt,minimum size=1mm,label=above:$M_0$] {};

\draw[->,NavyBlue,thick,vec] (45:3) --++ (-45:1.5) node[black,right] {$v_\text{к}$};
\node at (45:3) [circle,fill=,inner sep=0pt,minimum size=1mm,label={[inner xsep=0]above right:$M_1$}] {};

\draw[->,NavyBlue,thick,vec] (0:3) --++ (-90:1.5) node[black,right] {$v_\text{к}$};
\node at (0:3) [circle,fill=,inner sep=0pt,minimum size=1mm,label=right:$M_2$] {};

%omega 
\draw[->,thick,vec] (-3.2,0) arc [start angle=180, end angle=160, radius=3.2] node[left] {$\omega$};



    \end{tikzpicture}%
\caption{Движение по окружности}
\label{fig:p21}
\end{figure}

Муравей (не перемещаясь относительно колеса) последовательно проходит точки~$M_0$, $M_1$ и~$M_2$, лежащие на окружности радиусом~$R$. Тело имеет скорость~$v_\text{к}$, называемую \textbf{касательной скоростью} (или \emph{линейной скоростью}). В~данном движении модуль этой скорости сохраняется. 

Вращение обычно происходит \emph{по или против часовой стрелки}. Такое движение, после которого точка оказывается в начальном положении, есть \textbf{оборот}.

\textbf{Период} ($T$ [с])~--- это время, за которое тело совершает 1~\emph{оборот}. Можно видеть, что тело на рис.\,\ref{fig:p21} вернется в точку~$M_0$ через время $T=8$~с. Вот связи периода со скоростью и числом оборотов~$N$:
\begin{equation}
    T=\dfrac{2\pi R}{v_\text{к}}=\dfrac{\Delta t}{N}.
\end{equation}

\textbf{Частота} ($\nu$ [Гц])~--- это характеристика вращения, показывающая сколько оборотов совершает точка за одну секунду:
\begin{equation}
    \nu=\dfrac{1}{T}.
\end{equation}

\textbf{Угловая скорость} $\left(\omega\ \left[\dfrac{\text{рад}}{\text{с}}\right]\right)$~--- это характеристика вращения, показывающая на какой угол <<поворачивается>> \strut точка за одну~секунду. 
Скорость~$\omega$ как~бы указывает направление вращения точки в виде изогнутой стрелки (на рис.\,\ref{fig:p21} изображена черным цветом). Формула угловой скорости:
\begin{equation}
    \omega=\dfrac{\varphi}{\Delta t}=2\pi\nu,
\end{equation}
где $\varphi$~--- угол, на который <<повернулась>> точка.

Тело на рис.\,\ref{fig:p21} имеет ускорение~$a_\text{ц}$, называемое \textbf{центростремительным ускорением} (всегда направлено к центру окружности): 
\begin{equation}
    a_\text{ц}=\dfrac{v_\text{к}^2}{R}.
\end{equation}









 
\end{document}